\section{Алгоритм обратной задачи}
\label{sec:inverse-method}

\draft{Поток задан в 17 узлах на возрастах конфигураций и линейно
интерполируется между ними. Монотонность и границы обеспечиваются
параметрами $\eta_k\in[0,1]$:}
\begin{equation}
 q_0=q_{\min}+(q_{\max}-q_{\min})\eta_0,\qquad
 q_k=q_{\min}+(q_{k-1}-q_{\min})\eta_k,\quad k\geq1.
\end{equation}

\draft{Функционал~\eqref{eq:objective} минимизируется методом
POUNDERS~\cite{wild2017,petscPounders}, использующим модели отдельных невязок
без вычисления производных. Геометрия подготавливается однократно;
для каждого потока повторяются стационарный старт и прямой расчёт.
Промежуточные поля при подборе не записываются.}

\draft{Остановка: согласование всех наблюдений с погрешностью, внутренний
критерий оптимизатора или лимиты 10 итераций и 100 вычислений функционала.
Начальная история --- постоянные 75 мВт/м$^2$; перед запуском допускается
малый сдвиг параметров внутрь ограничений. После подбора отдельный прямой
расчёт сохраняет поля на исходной геометрии и найденный поток в XLSX.
Фактические истории потока и функционала показаны в разделе~\ref{sec:results}.}
